\section{Group or route? The same route flown alone}\label{sec:route}

A zone is 25~km wide. Inside it a convergence line, a ridge or a valley can let any pilot progress
straight, and can bring pilots together on it: the route would then explain both the group and the
straight progress. Does group time progress straighter than the same route flown alone?

\paragraph{Matching.} The candidates are the \RouteAlone{} paraglider flights that never had a contact
of any length. A site joins the takeoffs within 1~km across dates, and its flights are placed in one
east--north frame. A piece of group time and a candidate from the same site on another day match when
at least 80\% of the points of the piece, one every 30~s, lie within $\rho=500$~m of the candidate's
track, and the candidate flew them the same way. The candidate's segment, between its points nearest
to the first and to the last of them, is measured as the pieces of Section~\ref{sec:setup} if it lasts
at least 1250~s; each piece keeps at most 20 matches, the most covered first.

\paragraph{Comparison.} In analogy with Section~\ref{sec:zones}, the stratum is the piece: the
ensemble MSD of the matched pieces is compared with that of their segments reweighted to one piece
per stratum, Eq.~\eqref{eq:reweight}. Both sides are means of many MSDs before their logarithm. The
baseline compares the same pieces with the solo flights of their site-day and hour, so that the two
differences refer to the same group time. A second set keeps only the matches of the same season and
hour (month within one, local hour within two).

\begin{figure}
\centering
\includegraphics[width=\textwidth]{route_sample10000.pdf}
\caption{Left: $\Hloc(\tau)$ of the matched pieces of group time and of the same routes flown alone on
another day. Right: their difference over every lag, any date and same season and hour, with 5--95\%
bands.}
\label{fig:route}
\end{figure}

\begin{table}
\centering\footnotesize\setlength{\tabcolsep}{3pt}
\begin{tabular}{@{}l r c c c@{}}
\toprule
$\rho$, cover & pieces & route, any date & route, same season and hour & site-day and hour \\
\midrule
500~m, 80\% & 10\,000: 4\,933 & $+$0.018 [$+$0.014, $+$0.021] & $+$0.016 [$+$0.012, $+$0.021] & $+$0.060 [$+$0.055, $+$0.066] \\
500~m, 80\% & 2\,000: 958 & $+$0.021 [$+$0.013, $+$0.027] & $+$0.017 [$+$0.008, $+$0.024] & $+$0.061 [$+$0.051, $+$0.072] \\
300~m, 80\% & 2\,000: 675 & $+$0.022 [$+$0.012, $+$0.030] & $+$0.018 [$+$0.008, $+$0.029] & $+$0.066 [$+$0.053, $+$0.079] \\
1000~m, 80\% & 2\,000: 1\,284 & $+$0.014 [$+$0.010, $+$0.021] & $+$0.016 [$+$0.010, $+$0.023] & $+$0.057 [$+$0.049, $+$0.065] \\
500~m, 70\% & 2\,000: 1\,062 & $+$0.017 [$+$0.011, $+$0.025] & $+$0.016 [$+$0.009, $+$0.024] & $+$0.061 [$+$0.051, $+$0.070] \\
500~m, 90\% & 2\,000: 793 & $+$0.020 [$+$0.012, $+$0.028] & $+$0.015 [$+$0.006, $+$0.024] & $+$0.062 [$+$0.050, $+$0.073] \\
\bottomrule
\end{tabular}

\caption{$\Delta\Hloc(1000~\mathrm{s})$ of the matched pieces against the same route flown alone and
against the solo flights of their site-day and hour, by radius $\rho$ and cover. Pieces: drawn at
random, and matched. First row: the sample of the text.}
\label{tab:route}
\end{table}

\paragraph{Results.} Of a random sample of \RouteSample{} pieces, \RoutePieces{} have matches
(\RouteMatches{} segments, \RouteMedianMatches{} per piece in the median). Against the same route
flown alone $\Delta\Hloc(1000~\mathrm{s})$ is \RouteDHloc{} \RouteDHlocBand, and \RouteDHlocSeason{}
\RouteDHlocSeasonBand{} with the same season and hour (\RoutePiecesSeason{} pieces). The same pieces
against the solo flights of their site-day and hour give \RouteBaseDHloc{} \RouteBaseDHlocBand, more
than the \DHlocThousandSdh{} of all the group time, because these pieces lie on frequented routes.
The route therefore accounts for \RouteExplained\% of the long-lag gap (\RouteExplainedSeason\% with
the same season and hour), and the rest goes with flying together. Every radius and cover of
Table~\ref{tab:route} tells the same. At short lags group time is less straight than the same route
flown alone, $\Delta\Hloc(100~\mathrm{s})=$ \RouteDHlocHundred{} \RouteDHlocHundredBand, so the fitted
$H$ hardly differs (\RouteDH).

\paragraph{Limits.} The days differ: season and hour are matched, the weather of the day is not. The
pilots who never meet anyone may differ from those who fly in groups, and this test has no selection
contrast to measure it. Only the pieces on routes flown alone on other days enter, about half of them,
on a sample of the pieces.

\paragraph{What we keep.} About \RouteExplained\% of the straighter progress of groups over a quarter
of an hour comes from the routes they fly; flying together adds the rest.
